\documentclass[11pt]{article}
\usepackage[margin=1in]{geometry}
\usepackage[T1]{fontenc}
\usepackage[utf8]{inputenc}
\usepackage{amsmath,amssymb}
\usepackage{enumitem}
\setlength{\parindent}{0pt}
\setlength{\parskip}{0.75em}

\begin{document}

\begin{center}
{\Large \textbf{Referee Report}}\\[0.5em]
\textit{Computational resolution limits for super-resolution of positive sources}
\end{center}

This manuscript studies the super-resolution of positive point sources from noisy band-limited Fourier data.
More precisely, it analyzes computational resolution limits for two tasks:
detecting the correct number of sources and stably recovering their supports.
The authors obtain one-dimensional lower bounds in terms of the cutoff frequency,
the signal-to-noise ratio, and minimal separation, and then extend the theory to multi-dimensional settings.
The paper also includes numerical experiments exhibiting the predicted phase-transition behavior and discusses
algorithmic implications for sparsity-promoting reconstruction.

Super-resolution, line spectral estimation,
and deconvolution of point sources are central problems in inverse problems and
computational imaging, and positivity is a natural and practically relevant prior in many imaging applications.
The manuscript is therefore of interest to a substantial audience.
The paper addresses a computationally meaningful notion of resolution limit,
proves explicit quantitative bounds, and connects the theory to concrete reconstruction procedures and numerical phase transitions.

In my view, the strongest contribution of the paper is the lower bound for the \emph{number-detection}
resolution limit in the positive-source setting.
This seems to be the main new result relative to the authors' earlier works on general sources.
The lower bound present the same asymptotic order as the general sources (not necessarily positive) upper bound,
and is presented with explicit and fairly sharp constants.
Conceptually, the result is useful because it shows that positivity does \emph{not}
improve the asymptotic relationship between the super-resolution factor and the signal-to-noise ratio:
the exponents remain the same as in the general-source case.
At the same time, positivity does allow one to derive better constants.

The paper contains several additional useful contributions.
The support-recovery lower bound and instability constructions
help show that the general upper bounds are optimal in order for positive sources.
The multi-dimensional extension is natural and valuable.
I also found the observation in Section 4 interesting:
positivity might damage resolution, for two sources with different
phases the resolution limit can in fact be better than in the positive case.

My main reservation concerns the way the novelty is currently described relative to the previous paper [32].
After comparing the arguments, I do not think the lower-bound proof here should be presented as a wholly new methodology.
The proof architecture is largely the same as in [32]: an underdetermined Vandermonde-type system
is constructed, low-order moments are made to cancel, the coefficients are
controlled via inverse-Vandermonde/Lagrange-type formulas, and the Fourier
tail is estimated using factorial and Stirling-type bounds. What is new and
important here is the additional analysis needed to enforce \emph{positivity},
especially the sign analysis that allows the null vector to be split into two positive measures,
together with some extra combinatorial estimates in the more delicate constructions.
In other words, I view the argument as a meaningful positivity-adapted refinement of the authors'
earlier lower-bound technique, rather than a completely different technique.

The manuscript is otherwise nicely written and generally clear.
The results are well organized, the main theorems are easy to locate,
and the numerical sections help illustrate the theory. For these reasons,
my recommendation is: Accept subject to minor revisions.

\textbf{Comments to the authors}

\begin{enumerate}[leftmargin=2em]
\item \textbf{Clarify the main novelty more explicitly.}
In my opinion, the introduction should say more clearly that the main new one-dimensional
contribution is the lower bound for the number-detection problem in the case of positive sources.
At present, this point is somewhat diluted by the broader presentation.
It would help readers if Section~1.2 stated explicitly which upper bounds are
imported from previous work and which lower bounds are new here.

\item \textbf{Clarify the relation to the earlier general-source lower bounds.}
It would be very helpful to state explicitly,
at least in one dimension, the corresponding lower bounds from [32] for the general (not necessarily positive) case,
and then compare them with the present positive-source bounds.
The asymptotic orders are the same, whereas the constants appear to improve
in the positive setting; in one dimension, the gain seems to be noticeable
(roughly a factor of 4--6 for number detection and 7--8 for support recovery).
This comparison would make the contribution of the present paper much easier to assess.

\item \textbf{Moderate the description of the proof method.} I suggest revising the language around ``our new technique''
and instead describing the proof more precisely as an adaptation of the
lower-bound construction from [32] to the positive-source setting.
The main added ingredient, as I understand it, is the sign analysis ensuring positivity of the constructed measures.
Framing the contribution this way would be both accurate and still favorable to the paper.

\item \textbf{Express the main bounds also in SNR/SRF form.} Several results would be easier to compare with previous
work if they were rewritten in terms of the super-resolution factor and the signal-to-noise ratio.
For example, the number-detection threshold may be summarized in the form
\[
\mathrm{SNR} \gtrsim (C\,\mathrm{SRF})^{2n-2},
\]
and the support-recovery threshold in the form
\[
\mathrm{SNR} \gtrsim (C\,\mathrm{SRF})^{2n-1}.
\]
A short remark of this kind would improve readability and emphasize the scale-invariant content of the results.

\item \textbf{Please qualify Remark~2.2.} As written, the statement that the present estimates ``improve''
those of [32] could be read as suggesting an improved asymptotic rate.
My understanding is that the improvement is at the level of constants (and of adapting the theory to positive sources),
not at the level of the exponents. I suggest making this distinction explicit.

\item \textbf{A useful practical corollary could be stated after Theorem~2.3.}
By combining the separation condition with the pointwise localization estimate,
one can derive a simple bound showing that the recovered locations remain within a
fixed fraction of the minimum separation distance.
Even a brief remark of this type would help readers interpret the meaning of the theorem in concrete terms.

\item \textbf{A few small exposition improvements would help.}
\begin{enumerate}[label=(\alph*), leftmargin=2em]
\item A short note clarifying that the results are not essentially tied to the specific interval $I(n,\Omega)$,
but can be recentered and viewed on the real line, would be useful.
\item It would be good to standardize the notation for norms and absolute values throughout the paper.
\item Where appropriate, bibliographic entries and cross-references could be modernized slightly
(for example, by using hyperlinks in the final typeset version).
\end{enumerate}

\item \textbf{The discussion of previous literature on positivity could be sharpened slightly.}
I think the current discussion is already informative,
but a somewhat clearer separation between earlier results on the possibility
of super-resolution for positive sources, earlier algorithmic stability results,
and the present computational-resolution-limit viewpoint would improve the historical narrative.
In particular, it may be worth clarifying how this viewpoint relates to classical identifiability
or Prony-type arguments when sufficiently many measurements are available,
so that the reader can better distinguish the present stability-and-resolution-limit
perspective from earlier exact-recovery perspectives.
\end{enumerate}

Overall, I found the paper to be a useful and clearly written contribution.
The main new result is worthwhile, the constants are explicit and interesting,
and the manuscript is a good fit for the journal once the above points are addressed.

\end{document}